% Draft prepared from Jonas Seidel's MSc thesis (31 March 2026).
% Venue-independent working version; keep anonymous review mode until a venue is chosen.
\documentclass[acmsmall,screen,review,anonymous,nonacm]{acmart}

\usepackage[T1]{fontenc}
\usepackage{graphicx}
\usepackage[dvipsnames]{xcolor}
\usepackage{amsmath}
%\usepackage{amssymb}
\usepackage{stmaryrd}
%\usepackage{amsthm}
\usepackage{mathtools}
\DeclarePairedDelimiter\abs{\lvert}{\rvert}
\usepackage{mdframed}
\usepackage{booktabs}
\usepackage{tikz}
\usepackage{csquotes}
\usetikzlibrary{patterns,shapes,positioning,automata}
\usepackage[altpo]{backnaur}
\usepackage[shortlabels]{enumitem}
\usepackage{multicol}
\usepackage{adjustbox}
\usepackage{wrapfig}
\usepackage{nicefrac}
\usepackage{makecell}
\usepackage{siunitx}
\usepackage{colortbl}
\usepackage{listings}
\usepackage{cleveref}
\usepackage{thmtools}
\usepackage{thm-restate}
\usepackage{bussproofs}
\usepackage{xspace}
\usepackage{pifont}
%,disable
\usepackage[textsize=scriptsize,bordercolor=white]{todonotes}
\usepackage{setspace}
\setlength{\marginparwidth}{40mm} % make todo notes wider
%\setlength{\marginparsep}{2pt} % decrease space between margin notes and main text

\newcommand{\genericcomment}[3]{\todo[color=#1,fancyline]{\begin{spacing}{0.7}\textbf{#2:} #3\end{spacing}}}
\newcommand{\genericcommentinline}[3]{\todo[inline,color=#1,,author=#2]{#3}\xspace}
\newcommand{\kb}[1]{\genericcomment{violet!35}{KB}{#1}\xspace}
\newcommand{\kbin}[1]{\genericcommentinline{violet!35}{KB}{#1}\xspace}
\newcommand{\jpk}[1]{\genericcomment{red!35}{JPK}{#1}\xspace}
\newcommand{\jpkin}[1]{\genericcommentinline{red!35}{JPK}{#1}\xspace}
\newcommand{\dz}[1]{\genericcomment{Goldenrod!35}{DZ}{#1}\xspace}
\newcommand{\dzin}[1]{\genericcommentinline{Goldenrod!35}{DZ}{#1}\xspace}
\newcommand{\tw}[1]{\genericcomment{RoyalBlue!20}{TW}{#1}\xspace}
\newcommand{\twin}[1]{\genericcommentinline{RoyalBlue!20}{TW}{#1}\xspace}
% colors
\newcommand{\blue}[1]{\textcolor{blue}{#1}}
\newcommand{\red}[1]{\textcolor{Bittersweet}{#1}}
\newcommand{\gray}[1]{\textcolor{gray}{#1}}

%Preliminaries
\newcommand{\weight}[1]{|#1|}
\newcommand{\nweight}[1]{|#1|_{\prsqbrackop\prsqbrackcl}}
\newcommand{\subdists}{{\Deltasub \Sigma}}
\newcommand{\subdist}[1]{{\Deltasub #1}}
\newcommand{\subhalfdists}{{\Deltasubhalf \Sigma}}
\newcommand{\dists}{{\Delta \Sigma}}
\newcommand{\onlybdists}{\langle \bool \rangle}%Not used everywhere in the document
\newcommand{\onlynbdists}{\langle \neg \bool \rangle}%Not used everywhere in the document
\newcommand{\Deltasub}{\Delta_{\leq 1}}
\newcommand{\Deltasubhalf}{\Delta_{\leq 0.5}}
\newcommand{\Deltaspec}[1]{\Delta_{#1}}%Delta Speciale. (rarely used)
\newcommand{\statevals}{\textnormal{E}}
\newcommand{\LFPstateval}{e}
\newcommand{\LFPset}{X}%For the reach transformer
\newcommand{\LFPhoset}{X}%For the pre transformer: higher order
\newcommand{\Mfdr}{\initSet_{\textsc{fdr}}}
\newcommand{\varsfdr}{\vars_{\textsc{fdr}}}
\newcommand{\Mfldr}{\initSet_{\textsc{fldr}}}
\newcommand{\varsfldr}{\vars_\textsc{fldr}}
\newcommand{\stratsfdr}{\Pi^{\epsilon}_{\textsc{fdr}}}
\newcommand{\modelst}{\models_t}
\newcommand{\modelsp}{\models_p}
\newcommand{\initDist}{\mu_0} %Not used everywhere in the document?
\newcommand{\initStateVal}{\mu_0} %Not used everywhere in the document?
\newcommand{\initSet}{M} 
\newcommand{\postSet}{M} %FOR THE PRE TRANSFORMER
\newcommand{\statevalSet}{M} %Only once used
\newcommand{\hSetOne}{M} %hoare Sets (=(sub)-distribution sets)
\newcommand{\hSetTwo}{N}
\newcommand{\hSetThree}{O}
\newcommand{\hSetFour}{P}
\newcommand{\triple}[3]{\{\hspace{-2.1pt}|#1|\hspace{-2.1pt}\}~#2~\{\hspace{-2.1pt}|#3|\hspace{-2.1pt}\}} % braces could be improved
\newcommand{\tripleGen}{\triple{\hSetOne}{\prog}{\hSetTwo}} % generic triple
\newcommand{\inv}{I}
\newcommand{\statesSet}[4]{S_{\progvar{#1} \rightarrow #2}^{\progvar{#3} \rightarrow #4}}%WIRD auch für nur 2 Argumente Ohne Makro im Dokument verwendet
\newcommand{\marg}[2]{\mathbb{M}_{#1}(#2)}
\newcommand{\fixation}[3]{\mathbb{F}_{#1 \rightarrow #2 \sim #3}}
\newcommand{\unif}[3]{\mathsf{unif}_{#1}(#2,\dots,#3)}
\newcommand{\unifeps}[3]{\mathsf{unif}_{#1}^{\pm\epsilon}(#2,\dots,#3)}
\newcommand{\pro}[3]{\textnormal{Pr}_{#1}(#2 \mapsto #3)}
\newcommand{\prosim}[4]{\textnormal{Pr}_{#1}[#2 #3 #4]}
\newcommand{\diff}[1]{\textnormal{diff}(#1)}
\newcommand{\diffn}[2]{\textnormal{diff}_{#1}^{#2}}
\newcommand{\pr}[2]{\textnormal{Pr}_{#1}[#2]}
\newcommand{\minim}[2]{\min\{#1, #2\}}%Minimum
\newcommand{\dirac}[1]{\delta_{#1}} % Dirac distribution over state #1


%Conventional Letters
\newcommand{\prog}{C}
\newcommand{\loopbody}{B}
\newcommand{\progZero}{\prog_0}
\newcommand{\progOne}{\prog_1}
\newcommand{\progTwo}{\prog_2}
\newcommand{\progThree}{\prog_3}
\newcommand{\progFour}{\prog_4}
\newcommand{\arith}{E}
\newcommand{\bool}{\varphi}
\newcommand{\probability}{p}
\newcommand{\textvar}[1]{\normalfont #1}%Program variables IN TEXTS. That means when either only "variable v is..." occurs in text or VERY short math thins like "(v,c) is..." In longer math equations or "\sigma(v)..." always progvar macro is used. In lstlisting outside of math mode annotations, nothing is used.
\newcommand{\progvar}[1]{#1}%Program variables
\newcommand{\numb}[1]{\mathbf{#1}}%VALUES of IDENTICAL NAMES prog vars
\newcommand{\supp}{\textnormal{supp}}
\newcommand{\pgcl}{\textnormal{pGCL}}
\newcommand{\loops}{\textnormal{Loop}}
\newcommand{\pgclwhile}{\abs{\pgcl}_{\textnormal{while}}}
\newcommand{\progwhile}[1]{\abs{#1}_{\textnormal{while}}}
\newcommand{\conf}{\textnormal{Conf}}

%Einfache Namen
\newcommand{\vars}{\mathit{Var}}
\newcommand{\false}{\mathsf{f}}
\newcommand{\true}{\mathsf{t}}
\newcommand{\bools}{\mathbb{B}}
\newcommand{\Ifdr}{I^{\epsilon}_{\textsc{fdr}}}
\newcommand{\Ifldr}{I_{\textsc{fldr}}}
\newcommand{\Inim}{I_{\textsc{nim}}}
\newcommand{\progfldr}{\mathit{\prog_\textsc{fldr}}}
\newcommand{\progfdr}{\mathit{\prog_\textsc{fdr}}}
\newcommand{\prob}[2]{\textnormal{prob}_{#1}(#2)}
\newcommand{\probt}[2]{\textnormal{probt}_{#1}(#2)}
\newcommand{\prep}{\textnormal{prep}}
\newcommand{\lfp}{\mathsf{lfp\,}}
\newcommand{\nats}{\mathbb{N}}
\newcommand{\RgeqZero}{\mathbb{R}_{\geq 0}}
\newcommand{\RgeqZeroInf}{\mathbb{R}^{\infty}_{\geq 0}}
\newcommand{\RgZeroInf}{\mathbb{R}^{\infty}_{> 0}}
\newcommand{\supr}{\textnormal{sup}}
\newcommand{\infimum}{\textnormal{inf}}
\newcommand{\unmod}[1]{\textnormal{unmod}(#1)}

%Transformer
\newcommand{\reach}[1]{\mathsf{reach}\llbracket #1\rrbracket}
\newcommand{\post}[1]{\mathsf{sp}\llbracket #1\rrbracket}
\newcommand{\den}[1]{\mathsf{D}\llbracket #1\rrbracket}
\renewcommand{\wp}[1]{\textnormal{pre}\llbracket #1\rrbracket} %Should be the same, just old and new macro
\newcommand{\pre}[1]{\textnormal{pre}\llbracket #1\rrbracket} %Should be the same, just old and new macro


%n for normal: für wenn du die Namen der Transformer im Text benutzen möchtest
\newcommand{\reachn}{\mathsf{reach}}
\newcommand{\postn}{\mathsf{sp}}
\newcommand{\postcupn}{\mathsf{sp}^{\cup}}
\newcommand{\denn}{\mathsf{D}}
\newcommand{\wpn}{\textnormal{pre}} %Should be the same, just old and new macro
\newcommand{\pren}{\textnormal{pre}} %Should be the same, just old and new macro



%Inductive program structure
% UPPER CASE VERSION ARE SYNTAX FRAGMENTS TO BE USED IN ALIGN
\newcommand{\prfont}[1]{\textnormal{\texttt{#1}}}
\newcommand{\prspace}{\,}
\newcommand{\prparenop}{\prfont{(}}
\newcommand{\prparencl}{\prfont{)}}
\newcommand{\prsqbrackop}{\prfont{[}}
\newcommand{\prsqbrackcl}{\prfont{]}}
\newcommand{\prbraceop}{\prfont{\{}}
\newcommand{\prbracecl}{\prfont{\}}}
\newcommand{\ski}{\prfont{skip}}
\newcommand{\ass}[2]{#1\prspace\prfont{:=}\prspace#2}
\newcommand{\pchoice}[3]{\prbraceop#1\prbracecl \prspace \prsqbrackop#2\prsqbrackcl \prspace \prbraceop#3\prbracecl}
\newcommand{\nchoice}[2]{\{#1\} \prsqbrackop\prsqbrackcl \{#2\}}
\newcommand{\nchoicebox}{\prsqbrackop\prsqbrackcl}
\newcommand{\comp}[2]{#1\prspace\prfont{;}\prspace#2}
\newcommand{\SEMICLN}{\prspace\prfont{;}\prspace}
\newcommand{\ifelse}[3]{\prfont{if}\prspace\prparenop#1\prparencl\prspace\prbraceop#2\prbracecl\prspace\prfont{else}\prspace\prbraceop#3\prbracecl}
\newcommand{\ifelseskip}[2]{\prfont{if}\prspace\prparenop#1\prparencl\prspace\prbraceop#2\prbracecl}
\newcommand{\ELSE}{\prbracecl\prspace\prfont{else}\prspace\prbraceop}
\newcommand{\IF}[1]{\prfont{if}\prspace\prparenop#1\prparencl\prspace\prbraceop}
\newcommand{\while}[2]{\prfont{while}\prspace\prparenop#1\prparencl\prspace\prbraceop#2\prbracecl}
\newcommand{\WHILE}[1]{\prfont{while}\prspace\prparenop#1\prparencl\prspace\prbraceop}


%Preprocessing FLDR
\renewcommand{\H}[2]{H[#1,#2]}
\newcommand{\bound}[2]{\phi_{#1}(#2)}
\newcommand{\isThere}[2]{\Theta_{#1}(#2)}

%Shorthand Notation for Transformer Tables
%\newcommand{\postass}[2]{\overset{\rightharpoonup}{\ass{#1}{#2}}}
%\newcommand{\postass}[2]{\xrightharpoonup[\ass{#1}{#2}]{}}
%\newcommand{\preass}[2]{\xleftharpoonup[\ass{#1}{#2}]{}}
%\xrightharpoonup[below]{above}
\newcommand{\postass}[2]{\overset{\rightharpoonup}{\ass{#1}{#2}}}
\newcommand{\preass}[2]{\overset{\leftharpoonup}{\ass{#1}{#2}}}

\newcommand{\cl}[1]{\mathcal{M}_{#1}}%CLOSURE FOR BOOLS FOR PRE TRANSFORMER
\newcommand{\idSubDists}{\textnormal{id}_{2^{\subdists}}}

\newcommand{\clos}{c\ell}%CLOSURE FOR HOARE NEW WHILE RULES
\newcommand{\closOf}[1]{{\clos\,#1}}
\newcommand{\infW}[1]{\textnormal{Inf}_{\textnormal{w}}(#1)}%CLOSURE FOR HOARE NEW WHILE RULES
\newcommand{\closast}{\Delta_{\geq \infW{\inv}} \Sigma \cap \clos}
\newcommand{\closastinv}[1]{\Delta_{\geq \infW{#1}} \Sigma \cap \clos}%FOR SPECIFYING THE INV



%MDP Translation Stuff
\newcommand{\mdp}{\mathcal{M}}
\newcommand{\mdpTr}{\mdp_{\prog}}
%\newcommand{\strat}{\sigma} THIS IS NOW FOR PROGRAM STRATEGIES
\newcommand{\mdptuple}{(\states,\act,\prmdp)}
\newcommand{\states}{S}
\newcommand{\act}{{Act}}
\newcommand{\prmdp}{P}
\newcommand{\distsStates}{\Delta \states}
\newcommand{\subdistsStates}{\Deltasub \states}
\newcommand{\mdpDist}{\eta}
\newcommand{\mdpInitDist}{\eta_0}
\newcommand{\mdpInitDistTr}{\eta_{\initDist}}
\newcommand{\mdpInv}{\inv_{\mdp}}
\newcommand{\smallstep}[2]{\overset{#1.#2}{\longrightarrow}}
\newcommand{\bigstep}[1]{\overset{#1}{\rightsquigarrow}}
\newcommand{\bigstepn}{\rightsquigarrow}
\newcommand{\nextdist}[1]{#1}

%Markov Chains:
\newcommand{\markc}{\mathcal{M}}
\newcommand{\markctuple}{(\states,\prmdp)}
\newcommand{\markcTr}{\markc_{\prog}}

% from Tobi (intro)
\newcommand{\htrip}[3]{\blue{\{\hspace{-2.1pt}|#1|\hspace{-2.1pt}\}}~#2~\blue{\{\hspace{-2.1pt}|#3|\hspace{-2.1pt}\}}}
\newcommand{\intropre}{pre}
\newcommand{\intropost}{post}


\newcommand{\splitStateInline}[2]{\tikz[baseline=-2.75pt] \node [circle split,draw,rotate=90,inner sep=0.5pt] {\rotatebox{-90}{\tiny$#1$} \nodepart{lower} \rotatebox{-90}{\tiny$#2$}};}

% program annotation
\newcommand{\progAnno}[1]{\blue{\prfont{/\!\!/}~#1}}
\newcommand{\progAnnoNoSlashes}[1]{\blue{\textcolor{white}{\prfont{/\!\!/}}~#1}}

% space stuff
\newcommand{\eeq}{~{}={}~}
\newcommand{\iiff}{~{}\iff{}~}
\newcommand{\lland}{~{}\land{}~}



%Nondeterministic Choice Stuff
\newcommand{\lef}{l}
\newcommand{\rig}{r}
\newcommand{\progstrat}{\prog^{\strat}}
\newcommand{\denstrat}[1]{\mathsf{D}^{\strat}\llbracket #1\rrbracket}
\newcommand{\dens}[2]{\mathsf{D}^{#1}\llbracket #2\rrbracket}
\newcommand{\denhelper}[1]{\mathsf{H}^{\strat}\llbracket #1\rrbracket}
\newcommand{\denhelpernormal}{\mathsf{H}^{\strat}}
\newcommand{\denhelperf}[1]{\mathsf{F}^{\strat}\llbracket #1\rrbracket}
\newcommand{\denhelperfnormal}{\mathsf{F}^{\strat}}
\newcommand{\reachstrat}[1]{\mathsf{reach}^{\strat}\llbracket #1\rrbracket}
\newcommand{\poststrat}[1]{\mathsf{sp}^{\strat}\llbracket #1\rrbracket}
\newcommand{\posts}[2]{\mathsf{sp}^{#1}\llbracket #2\rrbracket}
\newcommand{\postcup}[1]{\mathsf{sp}^{\cup}\llbracket #1\rrbracket}
\newcommand{\postcap}[1]{\mathsf{sp}^{\cap}\llbracket #1\rrbracket}
\newcommand{\postall}[1]{\mathsf{sp}^{\forall}\llbracket #1\rrbracket}
\newcommand{\reachall}[1]{\mathsf{reach}^{\cup}\llbracket #1\rrbracket}
\newcommand{\strats}{\Pi}
\newcommand{\strat}{\pi}
\newcommand{\modelscup}{\models^{\cup}}
\newcommand{\postpi}[3]{\mathsf{sp}^{#1}_{#2}\llbracket #3\rrbracket}
\newcommand{\postmu}[3]{\mathsf{sp}^{#1}_{#2}\llbracket #3\rrbracket}

\newcommand{\stratsD}{\Pi}
\newcommand{\stratsMem}{\Pi_{\textnormal{Mem}}}
\newcommand{\stratsSigD}{\Pi_{\Sigma}}
\newcommand{\stratsN}{\Pi^{\Delta}}
\newcommand{\stratsSigN}{\Pi_{\Sigma}^{\Delta}}
\newcommand{\stratsSigNLC}{{}^{LC}\Pi_{\Sigma}^{\Delta}}
\newcommand{\stratsSigsD}{\Pi_{\Sigma^*}}
\newcommand{\stratsSigsN}{\Pi_{\Sigma^*}^{\Delta}}
\newcommand{\stratsCSigsD}{\Pi_{(\prog, \Sigma)^*}}
\newcommand{\stratsCSigsN}{\Pi_{(\prog, \Sigma)^*}^{\Delta}}
\newcommand{\stratsDeq}{{}^{=}\Pi}
\newcommand{\stratsSigDeq}{{}^{=}\Pi_{\Sigma}}
\newcommand{\stratsNeq}{{}^{=}\Pi^{\Delta}}
\newcommand{\stratsSigNeq}{{}^{=}\Pi_{\Sigma}^{\Delta}}
\newcommand{\stratsSigsDeq}{{}^{=}\Pi_{\Sigma^*}}
\newcommand{\stratsSigsNeq}{{}^{=}\Pi_{\Sigma^*}^{\Delta}}
\newcommand{\stratsCSigsDeq}{{}^{=}\Pi_{(\prog, \Sigma)^*}}
\newcommand{\stratsCSigsNeq}{{}^{=}\Pi_{(\prog, \Sigma)^*}^{\Delta}}
\newcommand{\stratsall}{\Pi^*}
\newcommand{\resprog}{\downharpoonright_{\sigma}}
\newcommand{\weaker}{\triangleright}
\newcommand{\loopc}{\textnormal{LC}}


%MDP Strategies
\newcommand{\tstrat}[1]{\mathfrak{S}_{#1}} %MDP strat that got translated from a (program) strat



%Nim Game
\newcommand{\Mnim}{\initSet_{\textsc{nim}}}
\newcommand{\prognim}{\mathit{\prog_\textsc{nim}}}
\newcommand{\stratnim}{\strat_{\textsc{nim}}}

%Tool names (used in related work)
\newcommand{\toolzar}{\textsf{Zar}\xspace}
\newcommand{\toolpsi}{\textsf{PSI}\xspace}

% wp shorthand in related work
\providecommand{\wpRelWork}{\wpn}



%New Commands
\newcommand{\Mhh}{\textnormal{Mhh}}
\newcommand{\idx}{\textnormal{idx}}
\newcommand{\proghh}{\textnormal{proghh}}
\newcommand{\Succ}{\textnormal{Succ}}

\newif\ifappendixversion
\appendixversiontrue

\begin{document}

\title{Program-Level Proof Rules for Almost-Sure Termination\\of Probabilistic Sampling Algorithms}

\author{Anonymous author(s)}

\begin{abstract}
Almost-sure termination (AST) is indispensable for turning partial-correctness arguments about randomized algorithms into total-correctness guarantees. Recent work provides a sound and relatively complete AST rule for probabilistic control-flow graphs (pCFGs), but applying it to source programs requires a translation and location-dependent certificates. We present a new proof rule directly at the program level. The rule interprets a loop body as a distribution transformer and separates two obligations: a non-negative supermartingale controls escape to unbounded regions, while a natural-valued variant decreases with uniformly positive probability and is bounded on every supermartingale sublevel set. Consequently, certificates describe complete loop iterations rather than individual control-flow edges. We prove the rule sound and relatively complete for probabilistic while-loops with loop-free bodies by relating it to the pCFG rule.

We first use a one-dimensional random walk to expose the rule’s core mechanism. We then compare pCFG-level and program-level proofs for the Fast Dice Roller, showing that both approaches apply, while the program-level proof reuses an existing distributional invariant and avoids location-indexed certificate engineering. Finally, we establish AST of the Fast Loaded Dice Roller and the Han--Hoshi sampler. These two case studies particularly highlight the practical advantages of the program-level rule: rejection and restart in the Fast Loaded Dice Roller are handled by a compact argument over complete tree-traversal iterations, while the dyadic refinement performed by Han--Hoshi admits an intuitive certificate based on the target-distribution boundaries contained in the current interval. In both cases, the program-level proofs are shorter and more direct than corresponding pCFG proofs, which would require certificates for every intermediate control location. Together with existing partial-correctness results, our AST proofs yield total-correctness guarantees for all three samplers.

\keywords{almost-sure termination, probabilistic programs, sampling algorithms, distribution transformers, loop invariants, supermartingales, Fast Dice Roller}
\end{abstract}

\maketitle

\section{Introduction}
\label{sec:intro}

Verifying an exact sampler requires both the intended output law and
almost-sure termination (AST). An argument about the relative probabilities
of terminating outputs alone does not exclude probability mass lost to
divergence. Termination is particularly subtle when rejection resets a
progress measure: that measure need not decrease on every outcome, or even
in expectation. Moreover, AST permits infinite expected runtime, as in the
symmetric random walk stopped at zero. These examples motivate
probabilistic progress conditions without uniform negative expected drift.

\paragraph{From control-flow edges to complete iterations.}
McIver et al.~\cite{10.1145/3158121} already give a source-level AST rule,
using one function with both supermartingale and progress obligations.
Majumdar and Sathiyanarayana~\cite{DBLP:journals/pacmpl/MajumdarS25}
establish a sound and relatively complete rule for probabilistic
control-flow graphs (pCFGs) with two separately chosen functions: a
natural-valued variant for progress and a non-negative supermartingale
that controls escape from regions where the variant is bounded.

Applying the pCFG rule requires certificates over control locations and
valuations. Assignments, tests, and random choices generate intermediate
obligations; strict descent on deterministic edges may require offsets
that merely account for control flow. Building on Seidel's program-level
development~\cite{seidel2026thesis}, we place the two-certificate argument
at loop heads, where each transition represents a complete iteration.
Users supply state functions and the body's post-distribution, without
annotating intermediate locations.

\paragraph{The iteration-level argument.}
At every enabled state, the variant decreases with uniformly positive
probability and the supermartingale does not increase in expectation.
The variant is bounded on each supermartingale sublevel set. This coupling
permits resets while controlling the regions in which progress is pursued.
Both functions vanish at termination; they may coincide but serve different
roles. Existing distributional invariants can retain output information,
although AST uses only their supports.

FDR illustrates the change in representation: doubling the sampling range
and recycling rejected values belong to one iteration. Its source
certificate captures their combined effect, while a pCFG certificate tracks
intermediate valuations (\Cref{fig:fdr-iteration-variants,fig:fdr-pcfg}).
The post-distribution must still be derived or supplied. Our comparison
concerns explicit annotations, not measured verification cost.

\paragraph{Contributions.}
\begin{enumerate}
\item \textbf{A two-certificate AST rule over distribution transformers.}
We formulate this rule for probabilistic while-loops with loop-free
bodies (\Cref{sec:rule}). We prove soundness and relative completeness
by translating certificates to and from pCFGs (\Cref{sec:metatheory}).
Completeness is relative to true arithmetic and is independent of the
restricted synthesis templates.
 \item \textbf{Sampler certificates and a representation comparison.}
 We instantiate the rule for FDR, the Fast Loaded Dice Roller (FLDR),
 and Han--Hoshi sampling
 (\Cref{sec:case-studies}), with an explicit pCFG comparison for FDR.
 FDR and FLDR reuse existing distributional invariants
 ~\cite{zilken2026verifyingsamplingalgorithmsdistributional}; Han--Hoshi
 supplies an interval-boundary certificate. These are applications of the
 rule to known samplers. The random walk illustrates its treatment of
 infinite expected runtime.
 \item \textbf{Certificate synthesis.}
 For integer-state inputs with two equiprobable outcomes, we derive
 sufficient constraints for the rule's premises and describe bounded
 counterexample-guided coefficient search (\Cref{sec:certificate-synthesis}).
 Invariants, iteration semantics, and any partition are supplied; the
 templates impose the stronger boundedness condition $U\leq V$.
 The shared implementation covers the random walk, symbolic FDR, an exact
 FLDR instance, and a structural FLDR abstraction. Fixed partitions extend
 the same search to piecewise affine certificates. SMT counterexample
 queries verify the resulting obligations over the entire supplied
 invariant domain; the connection between abstract FLDR and concrete
 preprocessing is proved separately.
\end{enumerate}

\paragraph{Scope.}
The rule considers fixed-probability discrete choices without
nondeterminism. Loop-free bodies have bounded path length; together with
fixed probabilities this supplies a positive lower bound on every
feasible body-path probability. The sampler loops above have this form,
with terminating preprocessing treated separately. Nested loops and
nondeterminism require extensions of the metatheory. Relative completeness
of the rule must be distinguished from the restricted synthesis search.

\section{Preliminaries}
\label{sec:prelims}

We fix a finite set $\vars$ of program variables and a countable state space $\Sigma$; a state $\sigma\in\Sigma$ maps variables to rational values. Arithmetic and Boolean expressions are interpreted as functions on states. For a Boolean expression $b$, let $\llbracket b\rrbracket=\{\sigma\mid \sigma\models b\}$ and let $[b]$ denote its Iverson bracket. Updating $x$ in $\sigma$ to $a(\sigma)$ is written $\sigma[x/a(\sigma)]$.

\subsection{Probabilistic programs and subdistributions}

We consider the probabilistic guarded-command language
\[
 C ::= \ski \mid \ass{x}{a} \mid \comp{C}{C} \mid \pchoice{C}{p}{C}
      \mid \ifelse{b}{C}{C} \mid \while{b}{C},
\]
where $p\in[0,1]\cap\mathbb{Q}$. A \emph{subdistribution} is a function $\mu:\Sigma\to\mathbb{R}_{\geq0}$ of weight $\weight{\mu}=\sum_{\sigma}\mu(\sigma)\leq1$. We write $\subdists$ for all subdistributions, $\supp(\mu)=\{\sigma\mid\mu(\sigma)>0\}$ for the support, and $\dirac{\sigma}$ for the Dirac distribution at $\sigma$. Missing mass represents nontermination.

Programs denote linear post-distribution transformers $\post{C}$. The clauses needed below are
\begin{align*}
 \post{\ski}(\mu) &= \mu,\\
 \post{\ass{x}{a}}(\mu)(\sigma)
   &= \sum_{\sigma'=\sigma[x/a(\sigma')]}\mu(\sigma'),\\
 \post{\pchoice{C_1}{p}{C_2}}(\mu)
   &= \post{C_1}(p\mu)+\post{C_2}((1-p)\mu),\\
 \post{\comp{C_1}{C_2}}(\mu)
   &= \post{C_2}(\post{C_1}(\mu)),\\
 \post{\ifelse{b}{C_1}{C_2}}(\mu)
   &= \post{C_1}([b]\mu)+\post{C_2}([\neg b]\mu).
\end{align*}
For a loop $C=\while{b}{B}$, let $F_\mu(e)=\mu+\post{B}([b]e)$. Its semantics is
\[
  \post{C}(\mu)=[\neg b]\cdot \lfp F_\mu.
\]
Equivalently, the Kleene approximants collect executions that leave the loop after finitely many iterations. Linearity gives
$\post{C}(\sum_i p_i\dirac{\sigma_i})=\sum_i p_i\post{C}(\dirac{\sigma_i})$, which lets us reduce one-iteration obligations to Dirac inputs.

\begin{definition}[Almost-sure termination]
A program $C$ is \emph{almost-surely terminating} on $M\subseteq\subdists$ if
$\weight{\post{C}(\mu)}=\weight{\mu}$ for every $\mu\in M$.
\end{definition}

Thus AST means that execution loses no probability mass. We do not require positive AST: the expected runtime may be infinite.

\subsection{Distributional invariants}

Our assertions are sets of distributions rather than predicates over individual states. This aligns the termination argument with existing partial-correctness proofs for sampling algorithms~\cite{zilken2026verifyingsamplingalgorithmsdistributional}.

\begin{definition}[Distributional invariant]
Let $C=\while{b}{B}$ and $M\subseteq\subdists$. A set $I\subseteq\subdists$ is a \emph{distributional invariant} for $C$ and $M$ if
\[
 M\subseteq I
 \quad\text{and}\quad
 [\neg b]\mu+\post{B}([b]\mu)\in I
 \quad\text{for all }\mu\in I.
\]
We abbreviate $\supp(I)=\bigcup_{\mu\in I}\supp(\mu)$.
\end{definition}

The invariant can simultaneously carry state-space bounds and genuinely distributional information such as equal probability of states in a row of a decision tree. The AST rule uses the support information, while retaining the stronger invariant permits direct composition with partial correctness.

\subsection{Supermartingales, variants, and sublevel sets}

A function $v:\supp(I)\to\mathbb{R}_{\geq0}$ is a one-iteration supermartingale for loop body $B$ if, for every enabled state $\sigma$,
\[
  \sum_{\sigma'\in\Sigma}\post{B}(\dirac{\sigma})(\sigma')v(\sigma')\leq v(\sigma).
\]
For $r\in\mathbb{R}_{\geq0}$, its sublevel set is
$v_{\leq r}=\{\sigma\in\supp(I)\mid v(\sigma)\leq r\}$.
A variant $u:\supp(I)\to\mathbb{N}$ need not decrease in expectation or on every outcome. We only require a uniformly positive chance to move to a smaller value. The role of $v$ is to ensure that executions cannot escape all regions in which $u$ has a finite upper bound.

This division of labour matters for rejection samplers. A rejected run may reset $u$ to a larger value; AST follows because progress remains possible with nonzero probability while $v$ controls the size of the relevant search space.

\section{Background: AST on Control-Flow Graphs}
\label{sec:background}

Our rule is derived from the pCFG characterization of Majumdar and Sathiyanarayana~\cite{DBLP:journals/pacmpl/MajumdarS25}. We recall only the ingredients needed to state the connection.

A probabilistic control-flow graph has configurations $(\ell,\sigma)$ consisting of a location and program state. Assignment and guard locations have a deterministic enabled successor; probabilistic locations choose a successor according to fixed edge probabilities. A terminal location $\ell_{\mathsf{out}}$ is absorbing. An inductive invariant $\mathit{Inv}$ contains the initial configurations and is closed under positive-probability transitions.

The pCFG rule asks for two certificates on $\mathit{Inv}$. A non-negative function $V$ is zero exactly at terminal configurations and is a supermartingale: it does not increase along deterministic transitions and does not increase in expectation at probabilistic locations. A natural-valued function $U$ is also zero exactly at termination. It strictly decreases along deterministic transitions; at probabilistic locations inside every sublevel set
\[
 V_{\le r}=\{(\ell,\sigma)\in\mathit{Inv}\mid V(\ell,\sigma)\le r\},
\]
there is a lower bound $\varepsilon_r>0$ on the probability of moving to a smaller $U$-value. Finally, $U$ is bounded on every $V_{\le r}$.

\begin{theorem}[pCFG AST rule, adapted from~\cite{DBLP:journals/pacmpl/MajumdarS25}]
\label{thm:pcfg-rule}
If a pCFG admits an inductive invariant and certificates $V,U$ satisfying the conditions above, then it reaches $\ell_{\mathsf{out}}$ almost surely. The rule is relatively complete with respect to the first-order theory used to express the certificates.
\end{theorem}

The soundness intuition has two layers. Inside a fixed sublevel set, $U$ ranges over a finite interval. Since every step offers a uniformly positive chance of descent, termination occurs almost surely unless the run leaves that set. The supermartingale $V$ makes indefinite escape through ever larger sublevel sets a probability-zero event. This accommodates null-recurrent examples such as the symmetric random walk.

The rule is semantically powerful but syntactically fine-grained. Suppose one source-level loop iteration consists of a guard, two assignments, a fair choice, and another conditional. A pCFG proof must assign $U$ and $V$ at each intermediate location. Constants and scaling factors are introduced merely to enforce descent across those internal edges. Our rule retains the same two-layer argument while contracting each finite execution of a loop-free body into one distribution-transformer step.

\paragraph{Translation assumption.}
The standard structural translation $\mathcal{T}[C]$ maps program syntax to a pCFG: assignments become update edges, conditionals become guarded edges, and probabilistic choice becomes a probabilistic location. Sequential composition connects exit and entry locations; a while-loop connects the body exit back to its guard. For a loop-free $B$, paths from its entry to exit correspond exactly to outcomes of $\post{B}(\dirac{\sigma})$:
\[
 \post{B}(\dirac{\sigma})(\sigma')
 = \Pr\bigl(\text{paths in }\mathcal{T}[B]
       \text{ from }(\ell_{\mathsf{in}},\sigma)
       \text{ to }(\ell_{\mathsf{out}},\sigma')\bigr).
\]
Consequently, $C$ is AST on $M$ exactly when $\mathcal{T}[C]$ reaches its exit almost surely from every state in $\supp(M)$. A full structural proof of this standard correspondence is summarized in \Cref{app:translation}; it is the bridge used in \Cref{sec:metatheory}.

\section{A Program-Level Rule for AST}
\label{sec:rule}

Consider a program $C=\while{b}{B}$ whose body $B$ is loop-free. Let $M$ be a set of initial subdistributions. We now state the main contribution.

\begin{theorem}[Program-level AST rule]
\label{thm:program-rule}
The loop $C$ is AST on $M$ if there exist:
\begin{enumerate}[label=(\arabic*)]
 \item a distributional invariant $I$ for $C$ and $M$;
 \item a function $v:\supp(I)\to\mathbb{R}_{\geq0}$ such that
   \begin{enumerate}[label=(\alph*)]
    \item $v(\sigma)=0$ iff $\sigma\in\supp(I)\cap\llbracket\neg b\rrbracket$;
    \item for every $\sigma\in\supp(I)\cap\llbracket b\rrbracket$,
    \[
      \sum_{\sigma'}\post{B}(\dirac{\sigma})(\sigma')v(\sigma')\leq v(\sigma);
    \]
   \end{enumerate}
 \item a function $u:\supp(I)\to\mathbb{N}$ such that
   \begin{enumerate}[label=(\alph*)]
    \item $u(\sigma)=0$ for every $\sigma\in\supp(I)\cap\llbracket\neg b\rrbracket$;
    \item there exists $\varepsilon>0$ for which, for every enabled state $\sigma\in\supp(I)\cap\llbracket b\rrbracket$;
    \[
      \sum_{\sigma':u(\sigma')<u(\sigma)}
        \post{B}(\dirac{\sigma})(\sigma')\geq\varepsilon;
      \tag{Decrease}
    \]
    \item for every $r\geq0$, the set
    $\{u(\sigma)\mid \sigma\in\supp(I),\ v(\sigma)\leq r\}$ is bounded.
   \end{enumerate}
\end{enumerate}
\end{theorem}

The conditions deliberately do not require $u$ to decrease in expectation. It may increase sharply on some outcomes, as happens when a rejection sampler restarts. The rule asks only for a fixed positive mass of outcomes that reduce $u$. Nor must $u$ be globally bounded. Its boundedness is local to sublevel sets of $v$, allowing null-recurrent computations.

\paragraph{Soundness proof intuition.}
Fix a sublevel set $v_{\le r}$ and let $K_r$ bound $u$ there. From every nonterminal state in this set, a complete loop iteration has probability at least $\varepsilon$ of decreasing $u$. Hence there is probability at least $\varepsilon^{K_r}$ of following enough decreasing iterations to reach zero, provided the run stays in the set. Repeating this argument rules out nontermination confined to $v_{\le r}$. Meanwhile, the supermartingale property bounds the probability of reaching arbitrarily large values of $v$; non-negative supermartingale convergence rules out escaping every finite sublevel set with positive probability. Thus all nonterminating behavior has measure zero.

\paragraph{Iteration-level abstraction.}
The essential difference to \Cref{thm:pcfg-rule} is the granularity of the decrease condition. For a body containing $m$ internal control-flow steps, the pCFG rule needs a value at each intermediate location and a descent-compatible branch at each probabilistic node. The new rule evaluates the final post-distribution once. Assignments, guards, and choices internal to $B$ are hidden behind $\post{B}$. This has three practical consequences:
\begin{enumerate}
 \item certificates are functions of program states, not location--state pairs;
 \item an existing distributional invariant can be reused directly; and
 \item algebraic simplification of $\post{B}(\dirac{\sigma})$ exposes exactly the outcomes relevant to progress.
\end{enumerate}

\paragraph{Uniform versus sublevel-dependent progress.}
The pCFG rule permits an $\varepsilon_r$ for each sublevel set. We use a global $\varepsilon$ because source programs here contain only fixed rational probabilistic choices, and because both samplers use fair bits. A more general statement can replace $\varepsilon$ by $\varepsilon_r$ throughout without changing the proof architecture.

\paragraph{Modularity.}
Initial deterministic code preceding the loop is handled by transforming $M$ into the set of loop-entry distributions. For a sampler with preprocessing, we regard the resulting finite data structure as read-only parameters of the sampling loop. The preprocessing itself terminates by ordinary bounded-loop reasoning. This isolates the probabilistic AST argument to the loop where it is needed.

\section{Sec 5}
\label{sec:5}

\section{Sec 6}
\label{sec:6}

\section{Sec 7}
\label{sec:7}

\section{Related Work}
\label{sec:relwork}

\paragraph{Proof rules for almost-sure termination.}
The classical approach to AST verifies a \emph{variant function} that decreases with a uniformly positive probability, going back to McIver and Morgan~\cite{DBLP:series/mcs/McIverM05}; their rule is complete only when the reachable state space is finite. Supermartingale-based rules relax this requirement. McIver et al.~\cite{10.1145/3158121} give a rule in which a variant, itself required to be a supermartingale, decreases with a state-dependent probability and satisfies an additional progress condition; this rule also handles demonic nondeterminism and extends the class of provably AST programs beyond what bounded-variant arguments allow. Majumdar and Sathiyanarayana~\cite{DBLP:journals/pacmpl/MajumdarS25} recently gave the first proof rule for AST that is both sound and relatively complete, formulated over probabilistic control-flow graphs (pCFGs). They further show that the earlier rule of McIver et al.~\cite{10.1145/3158121} enjoys the same relative-completeness guarantee when read at the program level. Our rule is derived from the pCFG rule of~\cite{DBLP:journals/pacmpl/MajumdarS25} and inherits its soundness and relative completeness, but it is stated directly over distribution-transformer semantics and factors the supermartingale and variant obligations into two independent certificates rather than a single combined one. This factorization is what lets rejection loops reset the variant without disturbing the supermartingale argument, which is central to our Fast Loaded Dice Roller (FLDR) case study. Qualitative AST is the property we target throughout; positive AST, which additionally bounds the expected runtime, is a strictly stronger and differently characterized property that we do not address.

\paragraph{Operational and denotational models.}
pCFGs and Markov decision processes are the standard operational substrate for termination and reachability analysis of probabilistic and randomized systems~\cite{DBLP:books/wi/Puterman94,DBLP:books/sp/BaierK08}. Tools such as Storm~\cite{DBLP:journals/sttt/HenselJKQV22} and refined analyses of visiting times and stationary behavior~\cite{DBLP:conf/tacas/MertensKQW24,DBLP:journals/jcss/JungesK0W21} operate at this level. Viewing probabilistic systems as \emph{distribution transformers} rather than as single-state transition systems has recently been fruitful for both Markov chains and MDPs: Aghamov et al.~\cite{DBLP:journals/corr/abs-2406-15087,DBLP:conf/birthday/AghamovBKNOPV25} model-check Markov chains through their action on distributions, and Akshay et al.~\cite{DBLP:conf/cav/AkshayCMZ23,DBLP:conf/ijcai/0001CMZ24} synthesize invariants and certified policies for MDPs under distributional reach-avoidance objectives; distribution-based objectives for MDPs are also studied in~\cite{DBLP:conf/lics/AkshayGV18}. Our rule applies the same distribution-transformer perspective one level higher, directly to program syntax rather than to an underlying automaton, which avoids the location-indexed bookkeeping that a translation to pCFGs or MDPs otherwise imposes.

\paragraph{Denotational semantics and program logics for probabilistic programs.}
The denotational, post-distribution-transformer semantics we build on originates with Kozen~\cite{DBLP:conf/focs/Kozen79,DBLP:conf/stoc/Kozen83} and underlies the weakest-precondition and weakest pre-expectation calculi of McIver and Morgan~\cite{DBLP:series/mcs/McIverM05} and Kaminski~\cite{DBLP:phd/dnb/Kaminski19}, in the tradition of Dijkstra's predicate-transformer semantics~\cite{DBLP:journals/cacm/Dijkstra75,DBLP:books/daglib/0067387}. Hoare-style program logics for probabilistic programs were introduced by den Hartog and de Vink~\cite{DBLP:journals/ijfcs/HartogV02} and extended by the Ellora framework~\cite{DBLP:conf/esop/BartheEGGHS18}; more recent calculi include quantitative strongest-postcondition reasoning~\cite{DBLP:journals/pacmpl/ZhangK22}, its hyperproperty generalization~\cite{DBLP:journals/corr/abs-2404-05097}, and modal separation logic for conditional probability~\cite{DBLP:journals/pacmpl/Li0H23}. These logics target partial correctness or information flow rather than AST; our rule complements them by supplying the missing termination argument needed to lift a partial-correctness proof to total correctness.

\paragraph{Verifying exact sampling algorithms.}
Exact discrete samplers built from fair coin flips, following the classical constructions of von Neumann~\cite{vonNeumann1951} and Lumbroso's Fast Dice Roller (FDR)~\cite{DBLP:journals/corr/abs-1304-1916}, have received growing verification attention. Saad et al.\ introduce the Fast Loaded Dice Roller (FLDR) as a near-optimal sampler for arbitrary discrete distributions~\cite{DBLP:conf/aistats/SaadFRM20}; the original paper gives an informal termination argument, while our contribution supplies a proof in the program logic. Bagnall, Stewart, and Banerjee's Zar compiler~\cite{DBLP:journals/pacmpl/Bagnall0023,10.1145/3591220} formally verifies samplers compiled from probabilistic programs with loops and conditioning in Coq, targeting the random-bit model; generating-function techniques have likewise been used to reason about probabilistic programs and about samplers such as the fast dice roller~\cite{DBLP:conf/lopstr/KlinkenbergBKKM20,DBLP:conf/cav/ChenKKW22,haase2026generatingfunctionsmeetoccupation}. Exact and guaranteed-bound inference engines such as PSI~\cite{DBLP:conf/cav/GehrMV16,DBLP:conf/pldi/GehrSV20}, the approach of Beutner et al.~\cite{DBLP:conf/pldi/BeutnerOZ22}, and scalable exact inference for discrete programs~\cite{DBLP:journals/pacmpl/HoltzenBM20} address the complementary problem of computing or bounding the output distribution rather than proving termination. Closest to our work is Zilken and Batz et al.'s framework of distributional invariants~\cite{zilken2026verifyingsamplingalgorithmsdistributional}, which proves partial correctness of FDR-style samplers by tracking distributional invariants under the same post-distribution-transformer semantics we use. That work establishes that FDR's output is uniformly distributed \emph{conditioned on termination}; our AST proofs for FDR and FLDR are the missing ingredient for total correctness, and our rule is deliberately phrased so that its distributional invariant obligation can reuse the one already constructed for partial correctness.

\paragraph{Nondeterminism.}
Several lines of work verify probabilistic programs or Markov models in the presence of demonic or angelic nondeterminism~\cite{10.1145/3158121,DBLP:journals/pacmpl/BatzBKW24,Wang2018ADS,Zilberstein2025,10.1145/3743131,Zilberstein_2025}. Our proof rule targets probabilistic while-loops without nondeterministic choice; extending the distribution-transformer formulation to resolve nondeterminism, e.g., by adapting strategy-based semantics as in~\cite{DBLP:journals/pacmpl/BatzBKW24}, is an interesting direction we do not pursue here.

\section{Conclusion and Future Work}
\label{sec:conclusion}

\paragraph{Achieved results.}
We presented a new proof rule for almost-sure termination of probabilistic while-loops with loop-free bodies, phrased entirely in distribution-transformer semantics. The rule separates two certificates: a non-negative supermartingale that bounds the probability of escaping to ever-larger sublevel sets, and a natural-valued variant that need only decrease with a uniformly positive probability and be bounded within each such sublevel set. We proved the rule sound and relatively complete by relating program-level certificates to the pCFG-level rule of Majumdar and Sathiyanarayana~\cite{DBLP:journals/pacmpl/MajumdarS25}: soundness follows by compiling our certificates into pCFG certificates, and relative completeness follows in the converse direction. On the one-dimensional random walk, the rule reproduces the textbook AST argument without any location-indexed bookkeeping. For the Fast Dice Roller, we gave AST proofs both at pCFG level and at program level; the comparison shows that the program-level proof is markedly shorter, reuses an existing distributional invariant, and avoids introducing artificial constants purely to enforce descent across intermediate control-flow edges. Finally, we proved almost-sure termination of the Fast Loaded Dice Roller, whose rejection-and-restart structure is captured by a single compact iteration-level argument: the supermartingale bounds the depth of the underlying decision tree, while the variant tracks the distance towards a non-rejecting leaf and resets, but does not grow unboundedly, on rejection. Combined with the existing partial-correctness result for such samplers~\cite{zilken2026verifyingsamplingalgorithmsdistributional}, our AST proofs complete the total-correctness argument for both FDR and FLDR. In particular, the FLDR proof closes the formal termination gap left by the original algorithm presentation.

\paragraph{Limitations.}
The rule as presented has two main restrictions. First, it applies to a single while-loop with a \emph{loop-free} body: the one-iteration distribution-transformer step that the rule reasons about is only well defined when the body itself contains no further loops. Nested loops can in principle be handled by applying the rule to the innermost loop first and treating its established AST as a black box when analyzing the enclosing loop, but we do not formalize or evaluate this modular composition here. Second, the underlying programming language does not include \emph{nondeterministic} (demonic or angelic) choice: every probabilistic branch has a fixed, syntactically given probability, and there is no construct for resolving nondeterminism via a scheduler. This mirrors the pCFG rule's treatment of nondeterminism-free models and keeps the connection between program syntax and one-step distribution-transformer semantics direct, but it excludes programs whose termination behavior depends on an adversarial or best-case choice of successor. As with the pCFG rule of~\cite{DBLP:journals/pacmpl/MajumdarS25}, we also restrict attention to programs over rational-valued variables, which avoids measure-theoretic subtleties but excludes continuous sampling primitives.

\paragraph{Future work.}
Lifting the loop-free-body restriction to a fully compositional treatment of nested and sequentially composed loops is the most immediate extension, and would let the rule be applied automatically along a program's syntax tree rather than loop by loop. Incorporating demonic or angelic nondeterminism, following strategy-based semantics for probabilistic programs with nondeterministic choice~\cite{DBLP:journals/pacmpl/BatzBKW24}, would connect the rule to settings where a scheduler resolves nondeterministic branching, as already supported by the pCFG rule of Majumdar and Sathiyanarayana~\cite{DBLP:journals/pacmpl/MajumdarS25}. It would also be worthwhile to study the precise relationship between our rule and the program-level rule of McIver et al.~\cite{10.1145/3158121} in more depth, in particular whether one can characterize classes of programs for which one of the two styles of certificate is systematically easier to construct than the other. Finally, the relative-completeness results now available for both pCFG-level and program-level AST rules~\cite{DBLP:journals/pacmpl/MajumdarS25,10.1145/3158121} suggest a concrete target for automation: since our rule decouples the supermartingale and variant obligations, each may be amenable to separate, specialized synthesis procedures, offering a path towards (semi-)automated AST proofs for sampling algorithms such as FDR and FLDR.


\ifappendixversion
\appendix
\section{Translation and AST Preservation}
\label{app:translation}

This appendix records the correspondence used by the metatheory. The structural translation $\mathcal{T}$ introduces entry and exit locations for every command. Assignment edges update the valuation, guarded edges encode conditionals and loop tests, and a probabilistic-choice location has outgoing probabilities $p$ and $1-p$. Sequential composition identifies the first command's exit with the second command's entry.

\begin{lemma}[Loop-free correspondence]
For loop-free $B$ and states $\sigma,\sigma'$,
\[
 \post{B}(\dirac{\sigma})(\sigma')
 =\Pr(\mathsf{Paths}_{\mathcal{T}[B]}((\ell_{\rm in},\sigma),(\ell_{\rm out},\sigma'))).
\]
\end{lemma}
\begin{proof}
By structural induction on $B$. Skip and assignment have one path. Sequential composition uses convolution over intermediate states. Conditional paths are partitioned by the guard. Probabilistic choice multiplies branch-path probabilities by $p$ and $1-p$, matching linearity of the post-transformer.
\end{proof}

For a loop, the $j$th Kleene approximant of its characteristic functional equals the probability distribution of pCFG runs that exit within at most $j$ visits to the guard. Taking suprema gives equality between post-distribution weight and pCFG reachability probability. Therefore $\weight{\post{C}(\dirac{\sigma})}=1$ iff the translated graph reaches its exit almost surely from $(\ell_{\rm in},\sigma)$. Linearity extends the equivalence to arbitrary initial subdistributions.

\section{Appendix to \Cref{sec:3}}
\label{app:3}
\section{Appendix to \Cref{sec:4}}
\label{app:4}
\section{Artifact and Proof-Development Notes}
\label{app:artifact}

The current manuscript is a venue-independent draft. Before submission, the following items should be completed: (i) mechanize or fully expand the translation-correctness proof; (ii) state the preprocessing properties of FLDR as explicit lemmas and connect them to its reference implementation; (iii) decide whether the generalized $\varepsilon_r$ version belongs in the main theorem; and (iv) evaluate proof size using a stable metric, for example number of certificate cases and generated verification conditions. The anonymous author block and review options should be replaced only after a venue has been selected.

\fi

\bibliographystyle{ACM-Reference-Format}
\bibliography{references}
\end{document}
